\documentclass[10pt,a4paper]{amsart}
\usepackage[margin=19mm]{geometry}
\usepackage{amsmath,amssymb,mathtools}
\usepackage[scr=rsfs,cal=euler]{mathalfa}
\usepackage{xfrac}
\usepackage[unicode,hidelinks,bookmarks=false]{hyperref}
\newcommand{\ie}{i.e.\ }
\newcommand{\cf}{cf.\ }
\newcommand{\resp}{respectively\ }
\newcommand{\Subsection}{\subsection}
\providecommand{\nobreakdash}{\nobreak\hbox{-}}
\setlength{\emergencystretch}{2em}
\multlinegap0pt
\usepackage{bm}
\usepackage{bbm}
\usepackage{dsfont}
\usepackage{xspace}
\usepackage{cancel}
\usepackage{mathtools}
\usepackage{amsthm}
\makeatletter
\@ifundefined{theorem}{\newtheorem{theorem}{Theorem}}{}
\@ifundefined{lemma}{\newtheorem{lemma}[theorem]{Lemma}}{}
\@ifundefined{proposition}{\newtheorem{proposition}[theorem]{Proposition}}{}
\makeatother
\newcommand{\R}{\mathbb{R}}
\title[Cartan's and Gauss's equations and rigidity theorems for iso]{Cartan's and Gauss's equations and rigidity theorems for isometric embeddings in low Sobolev regularity}
\author{Isaac Newell, Luc Nguyen}
\date{Source: arXiv:2607.02412v1 (CC BY 4.0)}
\begin{document}
\maketitle

\noindent\emph{Source and licence.} Cartan's and Gauss's equations and rigidity theorems for isometric embeddings in low Sobolev regularity. Version arXiv:2607.02412v1 (CC BY 4.0), distributed under the Creative Commons Attribution 4.0 International licence, as recorded in the arXiv licence metadata for this exact version and shown on the abstract page https://arxiv.org/abs/2607.02412v1. Full rights notice: licenses/arxiv-2607.02412-CC-BY-4.0.txt.\par\smallskip

\noindent\textit{Editorial scope.} Counted unit: the Lemma whose source label is \texttt{lemma:pullback\_commutation} in the article (source0.tex lines 425--430, source label \texttt{lemma:pullback\_commutation}), delivered together with the article's own complete original proof of it (source0.tex lines 431--464, one \texttt{proof} environment of 34 lines). No mathematics is written, repaired or re-derived: statement and proof are the article's own text line for line. Required context retained uncounted: lemma:comp\_diffeo\_wsp (lines 333--340); prop:WeakProducts (lines 379--399). Every definition, display and label of those lines is retained unchanged. Omitted from this package and not redistributed: 1--332, 341--378, 400--424, 465--1655. Nothing inside a retained excerpt is omitted or altered: every retained body is delivered exactly as the article's lines are written, so no omission receipt of any kind accompanies this package. Citations: the retained excerpts carry no citation command at all, so no bibliography entry of the article is transcribed and no reference is redistributed. Reference adaptation: none was needed and none was made; every reference inside the delivered unit resolves inside it, so no unresolved reference or citation is produced. Printed slips kept literal: no printed word, symbol, hypothesis or number of the counted statement or of its proof was altered. Layout and class: the article's own class is \texttt{article} with packages including graphicx, amssymb, amsthm, amsmath, enumerate, titlesec, dsfont, parskip, geometry, mathtools, bbm, hyperref; the wrapper is the lane's pinned \texttt{amsart} editorial wrapper at 10pt on A4, which restates the theorem-like environments the retained text needs and adds the article's own macro definitions that the retained text needs (\texttt{\textbackslash R}), with extra wrapper packages (bm, bbm, dsfont, xspace, cancel, mathtools). The delivered numbering is the unit's own. Assets: no retained span contains a figure environment, a graphics-inclusion command, a PDF-page inclusion, a table or a TikZ picture; no image file is copied into this package, the package declares no asset dependency and no image file is redistributed with it. Licence: version arXiv:2607.02412v1 of this article is distributed under the Creative Commons Attribution 4.0 International licence, as recorded in the arXiv licence metadata for this exact version and shown on the abstract page https://arxiv.org/abs/2607.02412v1. A full rights notice is delivered at licenses/arxiv-2607.02412-CC-BY-4.0.txt. Characters: every retained excerpt is pure ASCII, so the delivered engine, XeLaTeX with its default Latin Modern text font, reports no missing character.
    \begin{lemma} \label{lemma:comp_diffeo_wsp}
    Let $U, U' \subset \R^n$ be bounded Lipschitz domains in $\R^n$ and $\Psi:  U \to U'$ a $C^1$ diffeomorphism with inverse $\Phi$, such that $\Psi \in C^1(\overline U;\R^n)$ and $\Phi \in C^1(\overline U';\R^n)$. Let $0< s,\alpha < 1$, $1 \le p < \infty$. Then, $v \mapsto v \circ \Phi$ is a bounded linear map from $W^{s,p}(U)$ (respectively $c^{0,\alpha}(U)$ or $BV(U)$) into $W^{s,p}(U')$ (respectively $c^{0,\alpha}(U')$ or $BV(U')$). Furthermore, the following estimates hold with $L := \|D\Psi\|_{C^0(\overline U)}$ and $M := \|D\Phi\|_{C^0(\overline U')}$.
    \begin{align}
        [v\circ \Phi]_{W^{s,p}(U')} &\leq L^{\frac{2n}p} M^{\frac{n}{p}+s} [v]_{W^{s,p}(U)} \quad \forall v \in W^{s,p}(U), \label{eq:frac_Sob_diffeo}\\
        [v \circ \Phi]_{0,\alpha, U'| \epsilon} &\leq M^\alpha [v]_{0,\alpha,U|M\epsilon} \quad \forall v \in c^{0,\alpha}(U),\\
        \|v \circ \Phi\|_{BV(U')} &\leq C(n,U) L^{2n} \max \{M^{n+1}, 1\} \|v\|_{BV(U)} \label{eq:BV_est_for_v_circ_Phi} \quad \forall v \in BV(U).
    \end{align}
    \end{lemma}

    \begin{proposition} \label{prop:WeakProducts}
        Let $(M,g)$ be a compact $n$-dimensional Riemannian manifold and $U \subseteq M$ a smooth open subset. Assume that $1\leq k,l\leq n$ with $k+l \leq n$. Then the following hold.
        \begin{enumerate}
            \item[(i)] There exists a unique bounded bilinear map
            \begin{equation*}
                \mathcal{W} : CH^\frac12(U;\Lambda^kT^*M) \times H^{-\frac12}(U;\Lambda^l T^*M) \to CH^{-\frac12}(U;\Lambda^{k+l}T^*M)
            \end{equation*}
            such that for any $\alpha \in C^\infty(\overline U;\Lambda^k T^*M), \beta \in \mathcal{D}^l(U), \psi \in \mathcal{D}^{n-k-l}(U)$,
            \begin{equation*}
                \mathcal{W}(\alpha,\beta)[\psi] = \int_U \alpha \wedge \beta \wedge \psi.
            \end{equation*}
            \item[(ii)] There exists a unique bounded bilinear map
            \begin{equation*}
                \mathcal{I} : CH^\frac12(U;TM) \times H^{-\frac12}(U;\Lambda^k T^*M) \to CH^{-\frac12}(U;\Lambda^{k-1}T^*M)
            \end{equation*}
            such that for any $X \in C^\infty(\overline U; T\R^n)$, $\alpha \in \mathcal{D}^{k}(U)$, $\psi \in \mathcal{D}^{n-k+1}(U)$,
            \begin{equation*}
                \mathcal{I}(X,\alpha)[\psi] = \int_U (\iota_X \alpha) \wedge \psi.
            \end{equation*}
        \end{enumerate}
    \end{proposition}

    \begin{lemma} \label{lemma:pullback_commutation}
        Let $U,U',\Psi, \Phi$ be as above with $n=2$. Let $X \in CH^\frac12(U;T\R^2)$ and $\alpha \in H^\frac12(U;T^*\R^2)$. Then, with products defined in the sense of Proposition \ref{prop:WeakProducts}, the following equality holds in $CH^{-\frac12}(U';\Lambda^2 T^* \R^2)$:
        \begin{equation*}
            \Phi^*(\iota_{X}d\alpha) = \iota_{(\Psi_*X)} d(\Phi^*\alpha).
        \end{equation*}
    \end{lemma}

    \begin{proof}
        Let $\phi \in \mathcal{D}^1(U')$ be an arbitrary test one-form on $U'$. Then, for any sequences $\{X_\epsilon\} \subset C^\infty(\overline U;T\R^2)$ with $X_\epsilon \to X$ in $CH^\frac12(U;T\R^2)$ and $\{\alpha_\epsilon\} \subset C^\infty(\overline U;T^*\R^2)$ with $\alpha_\epsilon \to \alpha$ in $H^\frac12(U;T^*\R^2)$, we have by Proposition \ref{prop:WeakProducts} that
        \begin{align} 
                (\Phi^*(\iota_X d\alpha))[\phi] &= \iota_X d\alpha[\Psi^*\phi]\nonumber\\
                    &= \lim_{\epsilon \searrow 0} \int_U \big( \iota_{X_\epsilon}d\alpha_\epsilon\big) \wedge \Psi^*\phi\nonumber\\
                    &= \lim_{\epsilon \searrow 0} \int_U \Psi^*\Big( \Phi^*\big( \iota_{X_\epsilon}d\alpha_\epsilon\big) \wedge \phi\Big).
                    \label{eq:pullback_before_exchange}
        \end{align}
        Since, $X_\epsilon$ and $\alpha_\epsilon$ are smooth, the equality
        \begin{equation*}
            \Phi^*(\iota_{X_\epsilon} \alpha_\epsilon) = \iota_{\Psi_*X_\epsilon} (\Phi^*d\alpha_\epsilon),
        \end{equation*}
        holds classically (both sides are at least continuous). At this point we would like to compute $\Phi^*(d\alpha_\epsilon)$. For this, we claim  that
        \begin{equation} \label{eq:pullback_d_commute}
            \Phi^*(d\beta) = d(\Phi^*\beta) \quad \text{for any } \beta \in C^\infty(U;T^*\R^2)
        \end{equation}
        and both sides\footnote{A priori we only know that $\Phi^*\beta$ and $\Phi^*(d\beta)$ are $CH^\frac12$.} of the equation therefore belong to $CH^\frac12(U';\Lambda^2 T^* \R^2)$. Supposing this claim, we continue the proof of the lemma. By \eqref{eq:pullback_d_commute} we have 
        $\Phi^* d\alpha_\epsilon = d(\Phi^*\alpha_\epsilon)$ and thus $\Phi^*(\iota_{X_\epsilon} \alpha_\epsilon) = \iota_{\Psi_*X_\epsilon} d(\Phi^*\alpha_\epsilon)$. Inserting this into \eqref{eq:pullback_before_exchange} we have
        \begin{align*}
            (\Phi^*(\iota_X d\alpha))[\phi] &= \lim_{\epsilon \searrow 0} \int_U \Psi^* \big( \iota_{\Psi_*X_\epsilon}d(\Phi^*\alpha_\epsilon) \wedge \phi\big)\\
            &= \lim_{\epsilon \searrow 0} \int_{U'} \iota_{\Psi_*X_\epsilon}d(\Phi^*\alpha_\epsilon) \wedge \phi.
        \end{align*}
        Lemma \ref{lemma:comp_diffeo_wsp} implies that $\Psi_*X_\epsilon \to \Psi_*X$ in $CH^\frac12(U;T\R^2)$ and $\Phi^*\alpha_\epsilon \to \Phi^*\alpha$ in $H^\frac12(U;T^*\R^2)$. It follows from Proposition \ref{prop:WeakProducts} that $\iota_{\Psi_*X_\epsilon} d(\Phi^*\alpha_\epsilon) \to \iota_{\Phi_*X} d(\Phi^*\alpha)$ in $CH^{-\frac12}(U';T^*\R^2)$. Therefore the lemma will be proved once the claim \eqref{eq:pullback_d_commute} is proved.
        
        To prove the claim \eqref{eq:pullback_d_commute}, we write $\beta = \beta_i dx^i$, where $\beta_i$ are smooth functions. Then
        \begin{equation*}
            \Phi^*\beta = \beta_i(\Phi)d\Phi^i,
        \end{equation*}
        and, considering any smooth sequence $\Phi_\epsilon \to \Phi$ in $C^1(\overline U';\R^2)$, we have in the sense of distributions that 
        \begin{equation*}
            d(\Phi^*\beta) = \lim_{\epsilon\searrow 0} d\Big(\beta_i(\Phi)d\Phi^i_\epsilon\Big) = \lim_{\epsilon\searrow 0} \partial_j \beta_i(\Phi) d\Phi^j \wedge d\Phi^i_{\epsilon} = \Phi^*(d\beta),
        \end{equation*}
        as claimed. This concludes the proof.
    \end{proof}

\end{document}
